\documentclass[11pt]{article}
\usepackage[margin=1in]{geometry}

\usepackage{amsmath}
\usepackage{amsfonts}
\usepackage{amssymb}

\title{Corrected Solutions to Communication System Problems}
\author{Student at the University of South Carolina}
\date{\today}

\begin{document}

\maketitle

\section*{Problem 1: Binary ASK Scheme under AWGN}

Given: \( P_{T1} = 0.8 \), \( P_{T0} = 0.2 \), AWGN with \( N_0 = 0.2 \), and \( E_b = 1 \).

\subsection*{a. Optimum Linear Filter for Demodulation}
The optimum linear filter for demodulation in the presence of AWGN is the Matched Filter.

\subsection*{b. Optimum Threshold at the Detector}
The optimum threshold should be calculated considering the a priori probabilities and the noise variance. For equiprobable bits, the threshold would be zero. However, since \( P_{T1} \neq P_{T0} \), the threshold \( a \) is found by equating the likelihoods:
\[
a = \frac{E_b}{2} + \frac{N_0}{4} \ln\left(\frac{P_{T0}}{P_{T1}}\right)
\]
Substituting the given values:
\[
a = \frac{1}{2} + \frac{0.2}{4} \ln\left(\frac{0.2}{0.8}\right) \approx 0.43
\]

\subsection*{c. Bit-error Rate (BER)}
The bit-error rate \( P_b \) is calculated as:
\[
P_b = P_{T1} \cdot Q\left(\frac{a - \sqrt{E_b}}{\sqrt{N_0/2}}\right) + P_{T0} \cdot Q\left(\frac{\sqrt{E_b} - a}{\sqrt{N_0/2}}\right)
\]
Where \( Q(x) \) is the Q-function.

\section*{Problem 2: Digital Communication System}

\subsection*{Receiver Operation}
The continuous-time signal at the receiver, \( r_{Rx}(t) \), is the sum of the convolutions of the transmitted signals with the channel and added noise. The correct integral for the received signal \( r_{Rx}(t) \) is:
\[
r_{Rx}(t) = \int_{-\infty}^{\infty} r_{Tx}(\tau) p_{Rx}(t - \tau) d\tau + \int_{-\infty}^{\infty} w_n(\tau) p_{Rx}(t - \tau) d\tau
\]
where \( r_{Tx}(t) \) is the transmitted signal, \( w_n(t) \) is the AWGN, and \( p_{Rx}(t) \) is the impulse response of the receiver filter.

\end{document}
